\documentclass[10pt]{article}

\usepackage[utf8]{inputenc}

\usepackage[T1]{fontenc}

\usepackage{amsmath}

\usepackage{amsfonts}

\usepackage{amssymb}

\usepackage[version=4]{mhchem}

\usepackage{stmaryrd}



\begin{document}

этих значениях неизвестных, либо обе они ложны. Полная Э. т. $T$ конечной или счетной сигнатуры тогда и только тогда счетно категорична, когда для каждого $n$ существует конөчное число формул с $n$ свободными переменными $v_{1}, \ldots, v_{n}$ такое, что каждая формула соответствующей сигнатуры со свободными переменными $v_{1}, \ldots, v_{n}$ әквивалентна в $T$ одной из этих формул. Полная теория конечной или счетной сигнатуры, категоричная в одной несчетной мощности, категорична и во всякой другой несчетной мощности. Система $A$ сигнатуры $\Omega$ наз. ә л е м е н т а р н о й п о д с и с т ем 0 й системы $B$ той же сигнатуры, если $A$ является подсистемой системы $\boldsymbol{B}$ и для всякой формулы $\Phi\left(v_{1}\right.$, ..., $v_{n}$ ) Јогики предикатов 1-й ступени сигнатура $\Omega$ со свободными переменными $v_{1}, \ldots, v_{n}$ и всяких $a_{1}$, ..., $a_{n} \in A$ из истинности $\Phi\left(a_{1}, \ldots, a_{n}\right)$ в $\boldsymbol{A}$ следует истинность $\Phi\left(a_{1}, \ldots, a_{n}\right)$ в $B$. Э. т. $T$ наз. м о д е л ьн о по л н о й, если для любых ее моделей $A$ и $B$ из того, что $A$ является подсистемой системы $B$, следует, что $A$ является элементарной подсистемой $B$. Оказывается, что модельно полная теория, имеющая такую модель, к-рая изоморфоно вкладывается во всякую модель этої теории, является полной. Ун и в е р с а л ь н о э к в и в а л е н т ны м и наз. такие алгебраич. системы одной сигнатуры, на к-рых истинны одни и те же предваренные формулы, не содержащие кванторов существования. Модельно полная Ә. т., все модели к-рой универсально эквивалентны, является полной. Используя технику модельной полноты, мокно доказать полноту и разрешимость Э. т. вещественно замкнутых полей, в частности поля действительных чисел. Среди других Э. т., являющихся разрешимыми,- Э. т. сложения целых и натуральных чисел, Э. т. абелевых групп, поля $p$-адических чисел, конечных полей, полей вычетов, упорядотенных абелевых групп, булевых алгебр.



Теория неразрешимых Ә. т. была заложена А. Тарским (A. Tarski) в 40-е гг. 20 в., хотя неразрешимость Јогики предикатов первой ступени была доказана A. Чёрчем (A. Church) еще в 1936, а неразрешшимость арифметики натуральных чисел - Дж. Россером (J . Rosser) также в 1936. Говорят, что Ә. т. Th $(K)$ класса $K$ алгебраич. систем одной сигнатуры $\Omega$ н е о т д е лим а, если не существует рекурсивного множества формул, содержащего все тождественно истинные замкнутые формулы сигнатуры $\Omega$ и содержащегося в $\mathrm{Th}(K)$. Ә. т. класса $K_{1}$ систем сигнатуры $\left\langle P^{(2)}\right\rangle$, состоящей из одного двуместного предиката, наз. о т н о с и т ель ь о о п р е д е л и м о й в Ә. т. класса $K_{2}$ систем силнатуры $\Omega_{2}$, если существуют такие формулы $\Phi\left(v_{0}\right.$; $\left.u_{1}, \ldots, u_{s}\right),{ }^{2} \Psi\left(v_{1}, v_{2} ; u_{1}, \ldots, u_{s}\right)$ сигнатуры $\Omega_{2}$, что для каждой системы $A_{1}$ из $K_{1}$ можно найти такую систему $A_{2}$ из $K_{2}$ и такие элементы $b_{1}, \ldots, b_{s}$ в $A_{2}$, что множество $X=\left\{x \in A_{2} \mid \Phi\left(x ; b_{1}, \ldots, b_{s}\right)\right.$ истинно в $\left.A_{2}\right\}$ вместе с предикатом $P^{(2)}$, ошределенным на $X$ так, что $P^{(2)}(x, y)$ истинно тогда и только тогда, когда $\Psi(x, y$; $\left.b_{1}, \ldots, b_{s}\right)$ истинно в $A_{2}$, образует алгебраич. систему, изоморфную системе $A_{1}$. Это определение естественным образом распространяется и на теории классов $K_{1}$ произвольной сигнатуры. Если неотделимая Э. т. класса $K_{1}$ относительно определима в Ә. т. класса $K_{2}$, то Э. т. класса $K_{2}$ тоже неотделима. Это позволяет доказывать неотделимость Ә. т. многих классов алгебраич. систем. В качестве Th $\left(K_{1}\right)$ при этом удобно брать Э. т. всех конечных бинарных предикатов, Э. т. конечных симметричных предикатов и подобные Э. т. Неотделимые Э. т. неразрешимы. К неразрешимым Ә. т. относятся также Э. т. поля рациональных чисел, многих классов колец и полей. Важен результат А. И. Мальцева о неразрешимости Э. т. конечных групп.



JIum.: [1] E p ші о в Ю. Л.. [и др.] “Успехи матем. наук», 1965 , т. 20, в. 4, с. 37-108; [2] E p ш о в Ю. Л., Проблемы разрешимости и конструктивные модели, М., 1980.



Ю. Л. Ершов, М. А. Тайилин. ЭЛЕМЕНТАРНАЯ ТЕОРИЯ ЧИСЕЛ - раздеЛ чисељ теории, изучающий свойства чисел элементарными методами. Такие методы включают использование свойств делимости, различных форм аксиомы индукции и комбинаторные соображения. Иногда понятие элементарных методов расширяют за счет привлечения простейших элементов математич. анализа. Традиционно неәлементарными считают доказательства, в к-рых используются мнимые числа.



К Ә. т. ч. обычно относят задачи, возникающие в таких разделах теории чисел, как теория делимости, теория сравнений, теоретико-числовые функции, неопределенные уравнения, разбиения на слагаемые, аддитивные представления, приближения рациональными числами, цепные дроби. Нередко решение таких задач приводит к необходимости выходить за рамки әлементарных методов. Иногда вслед за отысканием неэлементарного решения какой-нибудь задачи находят и ее элементарное решение.



Задачи Ә. т. ч. имеют, как правило, многовековую исғорию и нередко стоят в истоках современных направлений теории чисел и алгебры.



Из сохранившихся клинописных таблиц древних вавилонян можно сделать вывод, что им не были чужды задачи разложения натуральных чисел на простые множители. В 5 в. до н. ә. пифагорейцы построили т. н. учение о четных и нечетных числах и обосновали предложение: произведение двух натуральных чисел четно тогда и только тогда, когда хотя бы один из сомножителей - четное число. Общая теория делимости, по существу, была построена Евклидом. В его "Началах" (3 в. до н. э.) вводится алгоритм отыскания наибольшего общего делителя двух целых чисел и на этой основе намечается обоснование основной теоремы арифметики целых чисел: о возможности и единственности разложения натурального числа в произведение простых сомножителей.



После того как К. Таусc (С. Gauss) в нач. 19 в. построил теорию делимости для целых комплексных чисел стало ясно, что изучение произвольного кольца нужно начинать с построения теории делимости в нем.



Все свойства целых чисел так или иначе связаны с простыми числами. Поэтому вопросы распределения простых чисел в натуральном ряду издавна интересовали ученых. Евклиду принадлежит первое доказательство бесконечности множества простых чисел. Только в сер. 19 в. П. Јl. Чебышев сделал следующий шаг в изучении функции $\pi(x)$ - числа простых тисел, не превосходяџих $x$. Ему удалось элементарными средствами доказать неравенства, из к-рых следовало, что



\[

0,92129 \frac{x}{\ln x}<\pi(x)<1,10555 \ldots \frac{x}{\ln x}

\]



для всех достаточно больших $x$. В действительности, $\pi(x) \sim \frac{x}{\ln x}$, но это удалось установить только в кон. 19 в. средствами комшлексного анализа. Долгое время считалось, что этот результат не может быть получен әлементарно. Однако А. Сельберг (А. Selberg, 1949) получил элементарное доказательство этой теоремы. П. Л. Чебышеву удалось также доказать, что для $n \geqslant 2$ интервал $(n, 2 n)$ содержит хотя бы одно простое число. Уточнение интервала, содержащего хотя бы одно простое число, требует более глубокого изучения поведения функции $\pi(x)$.



В 3 в. до н. э. был разработан метод Эратосфена решеma выделения простых чисел из множества натуральных чисел. В 1918 В. Брун (V. Brun) показал, что видоизменение этого метода может служить основой для изучения множества почти простых чисел. Им была получена Бруна пеорема о п р о с т ы х близне ц ах. Бруна решето является частным случаем общих решета ме-





\end{document}